\clearpage
\fancyhead[L]{\small GAUSS TO BODE.}
\hypertarget{native-438-part-lower}{}
\begin{center}
{\large GAUSS TO BODE.}

Astronomisches Jahrbuch für 1823 pp. 228..230. Berlin 1820.
\end{center}

\hfill Göttingen, 3 September 1820.

\par\indent As you requested, I am sending you the ephemeris for the course of \emph{Pallas} during its period of visibility in 1821. This time it was calculated by Mr. \textsc{von Staudt}, a young man of outstanding talent who is studying mathematics and astronomy here. The calculation was only finished today, and you wanted to receive it as soon as possible. I cannot therefore include the precise prediction of the opposition on 19 May 1821. It is still unfinished, and waiting for it would delay dispatch by one postal day.

I have given a detailed description of the \textsc{Reichenbach} meridian circle in our \textit{gelehrte Anzeigen}, so there is no need to repeat any of it here. I have probably made about 1200 observations with it since 21 February, but have so far reduced only some of them completely. This time I give you my observed solar declinations through 15 July (only some of the right ascensions observed at the same time with the transit telescope have been reduced). The position of the pole on the instrument has already been determined as precisely as possible from all the observations suited to the purpose. The observations have been corrected for refraction using \textsc{Bessel}’s table and for parallax, taking the mean horizontal parallax to be $=8''60$. No allowance has been made for flexure, however. My tests so far suggest that it is almost entirely imperceptible. If it does amount to perhaps \emph{a few tenths of a second}, only a long series of observations can determine that quantity.


\Needspace{17\baselineskip}
\SourceLeaf{439}
\begin{center}
\emph{Observed declinations of the Sun} 1820.

O upper limb \quad U lower limb
\end{center}

\begin{center}
\begin{tabular}{ll|r r r l}
\multicolumn{6}{c}{South} \\
Feb. & 22 & $10^\circ$ & $12'$ & $36''5$ & O \\
 & 27 & $8$ & $54$ & $23.0$ & U \\
 & 28 & $7$ & $59$ & $29.9$ & O \\
 & 29 & $8$ & $9$ & $11.3$ & U \\
Mar. & 5 & $6$ & $14$ & $26.5$ & U \\
 & 11 & $3$ & $53$ & $5.3$ & U \\
 & 14 & $2$ & $11$ & $0.0$ & O \\
 & 19 & $0$ & $44$ & $37.2$ & U \\
\multicolumn{6}{c}{North} \\
Apr. & 3 & $5$ & $38$ & $52.1$ & O \\
 & 4 & $6$ & $1$ & $48.2$ & O \\
\end{tabular}
\end{center}

\Needspace{15\baselineskip}
\begin{center}
\begin{tabular}{ll|r r r l}
\multicolumn{6}{c}{North} \\
Apr. & 12 & $9^\circ$ & $0'$ & $52''2$ & O \\
 & 13 & $8$ & $50$ & $43.1$ & U \\
 & 16 & $9$ & $55$ & $14.7$ & U \\
 & 18 & $11$ & $9$ & $15.3$ & O \\
 & 21 & $12$ & $11$ & $13.2$ & O \\
 & 22 & $11$ & $59$ & $26.2$ & U \\
 & 23 & $12$ & $19$ & $29.7$ & O \\
 & 24 & $13$ & $11$ & $9.9$ & O \\
 & 26 & $13$ & $50$ & $10.3$ & O \\
 & 27 & $13$ & $37$ & $30.7$ & U \\
May & 6 & $16$ & $19$ & $13.5$ & U \\
\end{tabular}
\end{center}

\Needspace{15\baselineskip}
\begin{center}
\begin{tabular}{ll|r r r l}
\multicolumn{6}{c}{North} \\
May & 11 & $18^\circ$ & $11'$ & $25''9$ & O \\
 & 19 & $20$ & $4$ & $23.5$ & O \\
 & 27 & $21$ & $4$ & $11.3$ & U \\
 & 31 & $21$ & $41$ & $11.8$ & U \\
June & 6 & $22$ & $25$ & $19.1$ & U \\
 & 20 & $23$ & $43$ & $27.5$ & O \\
 & 25 & $23$ & $50$ & $30.6$ & O \\
 & 27 & $23$ & $4$ & $58.1$ & U \\
July & 11 & $22$ & $23$ & $0.0$ & O \\
 & 12 & $22$ & $14$ & $51.1$ & O \\
 & 15 & $21$ & $16$ & $29.3$ & U \\
\end{tabular}
\end{center}

I also give you the results of my observations of 20 zenith stars, that is, their mean zenith distances reduced to the beginning of 1820. The reduction uses the same elements on which \textsc{Bessel}’s tables for the Pole Star are based. These are the same stars that Professor \textsc{Schumacher} observed with \textsc{Ramsden}’s zenith sector in Lauenburg in 1819 and in Skagen this year. From these observations the curvature of the different sections of the meridian arc from Skagen to Göttingen will be derived. I cannot give the curvatures yet, because I do not yet know the \emph{final results} of \textsc{Schumacher}’s observations. I add only that I observed each star \emph{at least} three times in one position of the circle and three times in the other. The reduction to the beginning of 1820 takes no account of the stars’ proper motion. This motion can, however, be derived with the greatest accuracy by comparing my determinations with \textsc{Brandley}’s, taking the altitude of the pole at our observatory to be $=51^\circ\,31'\,48''7$. The mean date of my observations falls in the first days of August.

\Needspace{27\baselineskip}
\begin{center}
\begin{tabular}{l|l|l}
Stars. & Mean Z. D. 1820. & Coll. error. \\
\hline
$\eta$ Cephei & $9^\circ\,36'\,44''28$ N. & $2'\,15''50$ \\
\emph{Piazzi XX.} 222. & $8\;\;17\;\;7.34$ & $16.65$ \\
$\iota$ Draconis & $8\;\;4\;\;11.09$ & $16.78$ \\
47 Draconis & $7\;\;38\;\;26.35$ & $17.19$ \\
Cephei 2 Hev. & $6\;\;50\;\;16.06$ & $16.96$ \\
48 Draconis & $6\;\;2\;\;55.95$ & $17.36$ \\
53 Draconis & $5\;\;1\;\;32.34$ & $17.06$ \\
33 Cygni & $4\;\;29\;\;25.71$ & $18.35$ \\
49 Draconis & $3\;\;52\;\;25.79$ & $17.18$ \\
46 Draconis & $3\;\;49\;\;49.51$ & $16.62$ \\
\emph{Piazzi XX.} 391. & $2\;\;17\;\;59.29$ & $17.64$ \\
51 Draconis & $1\;\;35\;\;38.56$ & $16.96$ \\
$\kappa$ Cygni & $1\;\;30\;\;37.18$ & $17.20$ \\
\emph{Piazzi XXI.} 32. & $1\;\;18\;\;6.93$ & $16.27$ \\
20 Cygni & $1\;\;0\;\;18.21$ & $17.54$ \\
$\iota$ Cygni & $0\;\;10\;\;46.28$ S. & $17.06$ \\
$c'$ Cygni Praec. & $1\;\;25\;\;6.27$ & $18.03$ \\
$c'$ Cygni Sequ. & $1\;\;25\;\;33.33$ & $16.37$ \\
$\theta$ Cygni & $1\;\;43\;\;15.97$ & $17.79$ \\
$\omega'$ Cygni & $2\;\;44\;\;22.59$ & $17.71$ \\
\multicolumn{2}{r|}{Mean collimation error . .} & $2'\,17''22$ \\
\end{tabular}
\end{center}

As you can see, the determinations of the collimation error differ from the mean by more than $1''$ for only two stars.
